% =============================================================================
% CONCEPTUAL ITEM: two recursive calls in W does NOT imply span drops to one.
% Targets the over-generalization of "the 2 becomes a 1" (see "Balanced?" part e).
% Counterexample verified in SML/NJ: sumTree (independent) vs sumAcc (dependent)
% compute the same value, same work-recurrence shape, different span.
% =============================================================================

\part[4]\hfill

  A student notices that for many tree functions, a work recurrence of
  $$W(n) = 2W(n/2) + O(1)$$
  goes with a span recurrence of
  $$S(n) = S(n/2) + O(1).$$
  They conclude: ``whenever the work recurrence has \textbf{two} recursive calls,
  the span recurrence has \textbf{one}.''

  Is this conclusion correct in general? If not, give a brief explanation and a
  concrete example of a function for which it fails.

  \begin{solutionorbox}[20em]
    No. Dropping from two recursive calls to one is only valid when the two
    recursive calls are \textbf{independent} --- neither one needs the other's
    result --- so that they may be evaluated in parallel. If the second call
    \textit{depends} on the first, they must run \textbf{sequentially}, and the
    span counts both.

    For example, compare these two (equivalent) ways to sum a tree:
    \begin{codeblock}
      fun sumTree Empty = 0
        | sumTree (Node (l, x, r)) =
            x + sumTree l + sumTree r

      fun sumAcc Empty acc = acc
        | sumAcc (Node (l, x, r)) acc =
            let
              val accL = sumAcc l acc
            in
              sumAcc r (accL + x)
            end
    \end{codeblock}

    Both have work $W(n) = 2W(n/2) + O(1) = O(n)$. But in \code{sumAcc}, the call
    \code{sumAcc r} takes \code{accL} as an argument, so it cannot begin until
    \code{sumAcc l} has finished. There is no parallelism to exploit, and
    $$S(n) = 2S(n/2) + O(1) = O(n),$$
    whereas \code{sumTree}'s independent calls do give
    $S(n) = S(n/2) + O(1) = O(\log n)$.

    So the shape of the \textbf{work} recurrence alone does not determine the
    span --- you must look at whether the recursive calls depend on one another.
  \end{solutionorbox}
